%============================================================
% CH1 · PATTERN 5 — FREQUENCY ARRAYS OVER SMALL VALUES: THEORY
%============================================================
\subsection{Pattern Theory}
\label{ssec:ch1-p5-theory}

\begin{patternbox}{Pattern 5 at a Glance}
When the values are small integers, use the \textbf{value as an index}: one pass fills a counting array \texttt{cnt[v]} = ``how many times does $v$ occur''. Questions about values --- duplicates, the most frequent value, how many are smaller, does any count repeat --- become questions about this short array.

\begin{center}
\begin{tikzpicture}[x=1cm, y=1cm]
    \dparray{a}{2.0}{$a$}{2/acell, 0/acell, 3/acell, 0/acell, 2/acell, 2/acell}
    \dpindex{v}{-0.15}{value $v$}{0}{3}
    \dparray{c}{0.5}{\texttt{cnt}}{2/compcell, 0/compcell, 3/compcell, 1/compcell}
    \dparray{l}{-0.95}{\texttt{less}}{0/basecell, 2/basecell, 2/basecell, 5/basecell, 6/basecell}
    \draw[dep] (a0.south) to[bend right=10] (c2.north);
    \draw[dep] (a2.south) to[bend left=10] (c3.north);
    \node[note, anchor=west, text width=5.2cm, align=left] at (5.6,0.4)
        {\texttt{cnt[v]}: occurrences of $v$ (one pass over $a$)\\[2pt]
        \texttt{less[v]} $= \texttt{cnt}[0] + \dots + \texttt{cnt}[v-1]$: elements smaller than $v$; the last cell is $n = 6$};
\end{tikzpicture}
\end{center}

\begin{recurrencebox}
\centering
\texttt{cnt[v - MIN]++} for every element \qquad\quad \texttt{less[v+1] = less[v] + cnt[v]}\\[3pt]
$\bigO{n + V}$ time, $\bigO{V}$ memory, where $V = \text{MAX} - \text{MIN} + 1$ is the size of the value range
\end{recurrencebox}
Works when $V$ is small (up to a few million). For values up to $10^9$ use a hash map (Chapter~\ref{ch:hashing}) or sort.
\end{patternbox}

%------------------------------------------------------------
\subsubsection*{From Pattern 4 to Pattern 5}

Pattern~4 computed indices from \emph{positions} ($2i$, $i \bmod n$) and, in Problem~\ref{ssec:ch1-build-permutation}, from \emph{values} ($i \to a[i]$), which worked because every value was a valid index. Pattern~5 makes that the whole method: a second array indexed by value. The single-pass aggregates of Pattern~1 (Section~\ref{sec:ch1-p1}) kept one number per array (a maximum, a count); a frequency array keeps one number \emph{per possible value}, still in one pass.

\begin{diagrambox}
\begin{tikzpicture}[x=1cm, y=1cm]
    \begin{scope}
        \node[font=\small\bfseries, text=deepblue, anchor=west] at (-0.6,1.65) {position $\to$ value};
        \dparray{p}{0.3}{}{7/acell, 5/aread, 9/acell}
        \aptr{p1}{i = 1}{successgreen}
        \node[note] at (1.35,-1.45) {$a[i]$: what is stored at position $i$};
    \end{scope}
    \draw[flow] (3.4,0.3) -- node[lbl, above]{swap roles} (5.0,0.3);
    \begin{scope}[xshift=5.6cm]
        \node[font=\small\bfseries, text=deepblue, anchor=west] at (-0.6,1.65) {value $\to$ count};
        \dpindex{q}{0.95}{}{0}{9}
        \dparray{r}{0.3}{}{0/compcell, 0/compcell, 0/compcell, 0/compcell, 0/compcell, 1/aread, 0/compcell, 1/compcell, 0/compcell, 1/compcell}
        \node[note] at (4.5,-0.5) {\texttt{cnt[5]}: how many positions hold 5};
    \end{scope}
\end{tikzpicture}
\figcap{A frequency array forgets \emph{where} values are and remembers only \emph{how often} each value occurs --- exactly what many questions need.}
\end{diagrambox}

%------------------------------------------------------------
\subsubsection*{When to recognise it --- and when not}

Use Pattern 5 when
\begin{itemize}[leftmargin=*]
    \item the constraints give a \textbf{small value range}: $0 \le a_i \le 100$, letters \texttt{a}--\texttt{z}, digits, $|a_i| \le 1000$, values $\le n$;
    \item the answer depends on values only through their \textbf{multiplicities}: duplicates, ``appears exactly $k$ times'', the mode, ``same multiset'' (E3.5), pairs with equal values or a given difference;
    \item you need, for every element, \textbf{how many elements are smaller} (or larger): prefix sums of counts (Problem~\ref{ssec:ch1-smaller-than-current});
    \item the question is about the counts themselves --- ``are all counts different?'': a frequency array \emph{of the counts} (Problem~\ref{ssec:ch1-unique-occurrences});
    \item you want to sort small integers in $\bigO{n + V}$ (counting sort).
\end{itemize}

\textbf{Do NOT use it} when the values are large or sparse ($a_i \le 10^9$): an array of $10^9$ counters does not fit in memory (4\,GB), and even $V = 10^7$ costs $\bigO{V}$ per test case. Then use a hash map (Chapter~\ref{ch:hashing}), sort and count runs, or compress the values first (Chapter~\ref{ch:sorting}). Do not use it when \emph{positions} matter (``the nearest equal element''): counts have forgotten them.

%------------------------------------------------------------
\subsubsection*{Tool 1 --- the counting array}

Allocate one counter per possible value, initialised to zero, and increment the counter of each element. If the values are $0..\text{MAX}$ the array needs $\text{MAX} + 1$ cells --- the classic off-by-one.

\begin{invariantbox}[Invariant of the counting loop]
After the first $i$ elements $a_0, \dots, a_{i-1}$ have been processed, $\texttt{cnt}[v] = |\{\, j < i : a_j = v \,\}|$ for every value $v$.

\textbf{Proof.} True for $i = 0$ (all zero). Processing $a_i$ increments exactly the counter of the value $a_i$, the only value whose number of occurrences among the first $i+1$ elements is larger than among the first $i$. \hfill$\blacksquare$

\textbf{Consequences.} $\sum_v \texttt{cnt}[v] = n$ (a free sanity check); the number of distinct values is the number of non-zero counters.
\end{invariantbox}

%------------------------------------------------------------
\subsubsection*{Tool 2 --- offsets for negative values}

Indices cannot be negative. If the values lie in $[\text{MIN}, \text{MAX}]$, store value $v$ in cell $v - \text{MIN}$; the array has $V = \text{MAX} - \text{MIN} + 1$ cells. For $|a_i| \le 1000$: \texttt{cnt[v + 1000]}, 2001 cells.

\begin{diagrambox}
\begin{tikzpicture}[x=1cm, y=1cm]
    \dpindex{vv}{1.25}{value $v$}{-3}{3}
    \dparray{o}{0.6}{\texttt{cnt}}{1/compcell, 0/compcell, 2/compcell, 1/compcell, 0/compcell, 0/compcell, 3/compcell}
    \dpindex{ii}{-0.05}{cell $v + 3$}{0}{6}
    \node[note, anchor=west, text width=4.4cm, align=left] at (6.7,0.6) {$a = [3,-1,0,3,-3,-1,3]$, $\text{MIN} = -3$: value $-1$ lives in cell 2, value 3 in cell 6};
\end{tikzpicture}
\figcap{An offset shifts the value range onto the indices $0..V-1$. Read the value of cell $k$ back as $k + \text{MIN}$.}
\end{diagrambox}

%------------------------------------------------------------
\subsubsection*{Tool 3 --- counts $\to$ prefix counts $\to$ ranks (a counting-sort preview)}

The number of elements smaller than $v$ is the sum of the counts of all values below $v$. One prefix-sum pass over \texttt{cnt} (prefix sums are the subject of Chapter~\ref{ch:prefix_sums}) answers it for every $v$:
\[
\texttt{less}[0] = 0, \qquad \texttt{less}[v+1] = \texttt{less}[v] + \texttt{cnt}[v].
\]
Then $\texttt{less}[v]$ elements are $< v$, $\texttt{less}[v+1]$ are $\le v$, and $n - \texttt{less}[v+1]$ are $> v$. In sorted order, the copies of $v$ occupy exactly the positions $[\texttt{less}[v], \texttt{less}[v+1])$. That is \textbf{counting sort}: walk over $a$ and put each element into the next free slot of its value's block. Walking in input order keeps equal elements in their original order (the sort is \emph{stable}), which matters when elements carry data. Chapter~\ref{ch:sorting} develops it further (radix sort).

\begin{diagrambox}
\begin{tikzpicture}[x=1cm, y=1cm]
    \dparray{s}{1.3}{$a$}{2/acell, 0/acell, 3/acell, 0/acell, 2/acell, 2/acell}
    \dparray{t}{-0.1}{sorted}{0/akeep, 0/akeep, 2/awin, 2/awin, 2/awin, 3/aread}
    \arange{t0}{t1}{$v = 0$: $[0, 2)$}{successgreen}
    \arange{t2}{t4}{$v = 2$: $[2, 5)$}{memoorange}
    \arange{t5}{t5}{$v=3$}{accentorange}
    \node[note, anchor=west, text width=4.8cm, align=left] at (5.8,0.6) {block of value $v$ = $[\texttt{less}[v], \texttt{less}[v+1])$ with $\texttt{less} = [0,2,2,5,6]$; value 1 has the empty block $[2,2)$};
\end{tikzpicture}
\figcap{Prefix counts are the start positions of the value blocks in the sorted array. The \emph{rank} of an element (number of strictly smaller elements) is the start of its block.}
\end{diagrambox}

%------------------------------------------------------------
\subsubsection*{Tool 4 --- frequency of frequencies}

Sometimes the question is about the counts: ``is every count different?'', ``how many values reach the maximal count?'', ``how many values occur exactly once?''. The counts are themselves small integers in $[1, n]$, so count them too: $\texttt{fof}[c]$ = number of distinct values that occur exactly $c$ times. Skip the zero counters (values that do not occur at all) --- forgetting to is the classic bug of Problem~\ref{ssec:ch1-unique-occurrences}.

\begin{diagrambox}
\begin{tikzpicture}[x=1cm, y=1cm]
    \dpindex{vv}{1.95}{value}{-3}{3}
    \dparray{c}{1.3}{\texttt{cnt}}{1/compcell, 0/emptycell, 2/compcell, 1/compcell, 0/emptycell, 0/emptycell, 3/compcell}
    \dpindex{ff}{-1.25}{count $c$}{1}{3}
    \dparray{f}{-0.6}{\texttt{fof}}{2/anscell, 1/basecell, 1/basecell}
    \draw[dep] (c0.south) to[bend right=15] (f0.north);
    \draw[dep] (c3.south) to[bend left=10] (f0.north);
    \draw[dep] (c2.south) to[bend left=5] (f1.north);
    \draw[dep] (c6.south) to[bend left=15] (f2.north);
    \node[note, anchor=west, text width=3.6cm, align=left] at (6.9,-0.6) {values $-3$ and $0$ both occur once, so $\texttt{fof}[1] = 2$: the counts are \emph{not} all different. Grey zeros are skipped.};
\end{tikzpicture}
\figcap{Counting the counts, on the array of Tool~2. Each distinct value contributes 1 to the cell of its count.}
\end{diagrambox}

%------------------------------------------------------------
\subsubsection*{Tool 5 --- when to switch to a hash map}

\begin{table}[H]
\centering\small
\renewcommand{\arraystretch}{1.2}
\begin{tabularx}{\linewidth}{@{}>{\raggedright\arraybackslash}p{3.6cm} >{\raggedright\arraybackslash}p{4.4cm} X@{}}
\toprule
\textbf{Value range} & \textbf{Tool} & \textbf{Cost} \\
\midrule
$V \le$ a few $10^6$ & counting array & $\bigO{n + V}$ time, $\bigO{V}$ memory, fastest in practice \\
$V$ huge, $n \le 10^6$ & \texttt{unordered\_map<int,int>} (Chapter~\ref{ch:hashing}) & $\bigO{n}$ expected; beware anti-hash tests on Codeforces \\
$V$ huge, order needed & sort, then count runs; or compress & $\bigO{n \log n}$, deterministic \\
\bottomrule
\end{tabularx}
\end{table}

\textbf{Many test cases.} On Codeforces the sum of $n$ over all tests is bounded, but the number of tests can be $10^4$. Re-allocating or clearing an array of size $V$ in \emph{every} test costs $10^4 \cdot V$: with $V = 10^6$ that is $10^{10}$. Size the counter by the current test ($n + 1$ cells when values are $\le n$) or undo only the cells you touched.

%------------------------------------------------------------
\subsubsection*{Tool 6 --- counting pairs by value}

A pair question that depends only on the two values is answered from the counts, without looking at pairs of positions:
\begin{itemize}[leftmargin=*]
    \item pairs $i < j$ with $a_i = a_j$: $\sum_v \binom{\texttt{cnt}[v]}{2} = \sum_v \texttt{cnt}[v](\texttt{cnt}[v]-1)/2$;
    \item pairs with $a_j - a_i = k > 0$ (in either order of positions): $\sum_v \texttt{cnt}[v]\cdot\texttt{cnt}[v+k]$ (E5.2);
    \item pairs with $a_i + a_j = s$: $\sum_{v < s - v} \texttt{cnt}[v]\cdot\texttt{cnt}[s-v]$, plus $\binom{\texttt{cnt}[s/2]}{2}$ when $s$ is even.
\end{itemize}
The cost drops from $\bigO{n^2}$ pairs to $\bigO{V}$ values. Characters are small values too: \texttt{cnt[ch - 'a']} counts lowercase letters in 26 cells.

\begin{diagrambox}
\begin{tikzpicture}[x=1cm, y=1cm]
    \dparray{e}{0.8}{$a$}{2/gvA, 0/gvB, 3/acell, 0/gvB, 2/gvA, 2/gvA}
    \dpindex{ev}{-0.05}{value}{0}{3}
    \dparray{ec}{-0.7}{\texttt{cnt}}{2/compcell, 0/compcell, 3/compcell, 1/compcell}
    \node[note, anchor=west, text width=6.2cm, align=left] at (5.9,0.1) {equal pairs: $\binom{2}{2} + \binom{0}{2} + \binom{3}{2} + \binom{1}{2} = 1 + 0 + 3 + 0 = 4$ --- the pair of red 0s and the three pairs among the blue 2s};
\end{tikzpicture}
\figcap{Counting equal pairs from the counts: each value contributes $\binom{\texttt{cnt}}{2}$, independently of where its copies are.}
\end{diagrambox}

%------------------------------------------------------------
\subsubsection*{Variations map}

\begin{diagrambox}
\begin{tikzpicture}[
    var/.style={rectangle, rounded corners=2pt, draw=patternborder, fill=patternbg,
                font=\scriptsize, align=center, text width=3.25cm, minimum height=1.25cm, inner sep=3pt}]
    \node[var, fill=answerfill, draw=accentorange, text width=5.2cm, font=\scriptsize\bfseries] (root) at (0,0)
        {Frequency arrays\\\texttt{cnt[v - MIN]++}, one pass};
    \node[var] (v1) at (-5.4,-2) {\textbf{ranks}\\prefix counts: how many are smaller (Problem~\ref{ssec:ch1-smaller-than-current})};
    \node[var] (v2) at (-1.8,-2) {\textbf{counts of counts}\\all counts distinct? (Problem~\ref{ssec:ch1-unique-occurrences})};
    \node[var] (v3) at (1.8,-2)  {\textbf{special counts}\\$\texttt{cnt}[v] = v$, shortfalls $\max(0, m - \texttt{cnt})$ (E5.1, E5.3)};
    \node[var] (v4) at (5.4,-2)  {\textbf{pairs by value}\\$\sum_v \texttt{cnt}[v]\cdot\texttt{cnt}[v+k]$ (E5.2)};
    \node[var] (v5) at (-3.6,-3.6) {\textbf{bounded counts}\\every count $\le 2$? (E5.4)};
    \node[var] (v6) at (0,-3.6)    {\textbf{capped buckets}\\values $> n$ collapse into bucket $n$ (E5.5)};
    \node[var] (v7) at (3.6,-3.6)  {\textbf{distinct vs.\ mode}\\number of distinct values and the largest count (E5.6)};
    \node[draw=patternborder, dashed, rounded corners=4pt, fit=(v1)(v4)(v5)(v7), inner sep=5pt] (grp) {};
    \draw[flow] (root.south) -- (grp.north);
\end{tikzpicture}
\figcap{Every box reads only the counting array (and maybe its prefix sums or its own counts). Hash-map counting of large values is Chapter~\ref{ch:hashing}.}
\end{diagrambox}

%------------------------------------------------------------
\subsubsection*{Code templates}

\begin{lstlisting}[style=cpp]
#include <bits/stdc++.h>
using namespace std;

const int MINV = -5, MAXV = 5;                 // value range of this example
const int V = MAXV - MINV + 1;                 // number of possible values

int main() {
    vector<int> a = {2, -1, 3, -1, 0, 2, 2, -5};
    int n = a.size();
    // (1) counting with an offset: value v lives in cnt[v - MINV]
    vector<int> cnt(V, 0);
    for (int x : a) cnt[x - MINV]++;
    // (2) prefix counts: less[v - MINV] = how many elements are < v
    vector<int> less(V + 1, 0);
    for (int k = 0; k < V; k++) less[k + 1] = less[k] + cnt[k];
    for (int x : a) cout << less[x - MINV] << " ";          // rank = number of smaller elements
    cout << "\n";
    // (3) counting sort (stable): each element goes to the next free slot of its value
    vector<int> pos(less.begin(), less.end() - 1), sorted(n);
    for (int x : a) sorted[pos[x - MINV]++] = x;
    for (int x : sorted) cout << x << " ";
    cout << "\n";
    // (4) frequency of frequencies: how many distinct values occur exactly c times
    vector<int> fof(n + 1, 0);
    for (int c : cnt) if (c > 0) fof[c]++;
    for (int c = 1; c <= n; c++) if (fof[c]) cout << c << ":" << fof[c] << " ";
    cout << "\n";
    // (5) the mode (most frequent value, smallest on ties)
    int best = 0;
    for (int k = 1; k < V; k++) if (cnt[k] > cnt[best]) best = k;
    cout << best + MINV << " x" << cnt[best] << "\n";
}
\end{lstlisting}

\begin{insightbox}[Check the templates by hand]
$a = [2,-1,3,-1,0,2,2,-5]$ sorted is $[-5,-1,-1,0,2,2,2,3]$. The ranks are the first positions of each value there: $2 \to 4$, $-1 \to 1$, $3 \to 7$, $0 \to 3$, $-5 \to 0$; so the first line is \texttt{4 1 7 1 3 4 4 0}. Counts: $-5, 0, 3$ once, $-1$ twice, $2$ three times, hence \texttt{1:3 2:1 3:1}, and the mode is $2$ with 3 copies. All four lines come from running the program.
\end{insightbox}

%------------------------------------------------------------
\subsubsection*{Common mistakes}

\begin{warningbox}[Mistake 1 --- one cell too few]
\begin{lstlisting}[style=snippet]
vector<int> cnt(100);       // WRONG for values 0..100: cnt[100] is out of bounds
vector<int> cnt(100 + 1);   // RIGHT: MAX + 1 cells (MAX - MIN + 1 with an offset)
\end{lstlisting}
\end{warningbox}

\begin{warningbox}[Mistake 2 --- a negative index]
\texttt{cnt[x]++} with $x = -3$ writes before the array: no error message, just corrupted memory. Every negative value range needs an offset.
\end{warningbox}

\begin{warningbox}[Mistake 3 --- counting the zero counts]
When scanning \texttt{cnt} for a property of ``the values that occur'', skip $\texttt{cnt}[v] = 0$. Otherwise ``are all counts different?'' fails as soon as two values are absent.
\end{warningbox}

\begin{warningbox}[Mistake 4 --- clearing a big array in every test case]
\texttt{vector<int> cnt(1000001)} inside a loop over $10^4$ test cases is $10^{10}$ operations. Size it by the test, or reset only the touched cells.
\end{warningbox}

\begin{warningbox}[Mistake 5 --- a counting array for huge values]
$a_i \le 10^9$ needs 4\,GB of counters. Read the constraints; switch to sorting or hashing.
\end{warningbox}

\begin{checklistbox}[``I'm stuck'' checklist for Pattern 5]
\begin{itemize}[leftmargin=*]
    \item[$\square$] What is the value range? Is $V$ small enough for an array? Do I need an offset?
    \item[$\square$] Does the answer depend only on how often each value occurs? Then positions can be forgotten.
    \item[$\square$] Do I need ``how many are smaller / larger''? Prefix sums of the counts.
    \item[$\square$] Is the question about the counts themselves? Count the counts, skipping zeros.
    \item[$\square$] Many test cases: is my counting array sized or reset per test?
    \item[$\square$] Values too large? Hash map, or sort and count runs.
\end{itemize}
\end{checklistbox}
